\documentclass{article}
\input{../../../lib.sty}

\title{\texttt{fn make\_stable\_truncate}}
\author{\MichaelShoemateAuthor}
\begin{document}
\maketitle


\contrib

Proves soundness of \rustdoc{transformations/make_stable_lazyframe/truncate/fn}{make\_stable\_truncate}
in \asOfCommit{mod.rs}{a105afe}.

\section{Hoare Triple}
\subsection*{Precondition}
\subsubsection*{Compiler-verified}
None

\subsubsection*{Caller-verified}
The nonempty matched truncation pipeline is produced by \rustdoc{transformations/make_stable_lazyframe/truncate/matching/fn}{match\_truncations}.

\subsection*{Pseudocode}
\label{sec:python-pseudocode}
\lstinputlisting[language=Python,firstline=2,escapechar=|]{./pseudocode/make_stable_truncate.py}

\subsection*{Postcondition}
\begin{theorem}
\validTransformation{\texttt{(input\_domain, input\_metric, matched)}}{\\\texttt{make\_stable\_truncate}}
\end{theorem}

\begin{proof}[Appropriate Output Domain]
    By line \ref{line:stable-keys},
    the identifier and all grouping keys must be stable row-by-row expressions.
    In particular, filtering other rows must not change the identifier of a retained row.
    The identifier is unchanged because it defines the input metric.
    On line \ref{line:check-rewrites}, per-group truncation rewrites only partition keys;
    number-of-groups truncation uses the same key extractor as matching to replace
    whole occurrences of the grouping expressions in the dense rank input.
    Legacy sort keys are checked for permutation equivariance against the preceding plan's schema.
    The rank structure and threshold are unchanged, and the corresponding bounds
    are expressed using these same execution keys, including the terminal group-by bound.
    Line \ref{line:allowed_keys_check} checks that filter grouping keys are retained by the final group-by.
    Filters remove rows without changing the schema, preserving margin upper
    bounds but potentially invalidating margin invariants.
    On line \ref{line:truncation}, if a terminal groupby is present, the shared
    group-by preparation checks grouping keys, rewrites and certifies
    the aggregates, infers the output schema, translates retained margins to
    materialized output columns, and discards their invariants, as established in
    the output-domain argument for
    \rustdoc{transformations/make_stable_lazyframe/group_by/fn}{make\_stable\_group\_by}.
    Otherwise, all margin invariants are discarded once.
    Hence every output is in the derived domain or raises a data-independent runtime exception.
\end{proof}

\begin{proof}[Stability guarantee]
    By line \ref{line:total-ids}, \texttt{total\_num\_ids} is the total number of ids an individual may contribute to a dataset:

    \begin{equation}
        d_{\mathrm{SymId}}(\texttt{function}(x), \texttt{function}(x')) \le \texttt{total\_num\_ids}.
    \end{equation}

    By the postcondition of \rustdoc{transformations/make_stable_lazyframe/truncate/matching/fn}{match\_truncations},
    for each \texttt{truncation\_bound} on line \ref{line:truncation-bound},
    \begin{align*}
        \max_{id} ||d_{\mathrm{Sym}}(\texttt{function}(x)_{id,g}, \texttt{function}(x')_{id,g})||_\infty &\le \texttt{truncation.per\_group}, \\
        \max_{id} ||d_{\mathrm{Sym}}(\texttt{function}(x)_{id,g}, \texttt{function}(x')_{id,g})||_0 &\le \texttt{truncation.num\_groups},
    \end{align*}
    where $g$ labels a group defined by the certified execution expressions in \texttt{truncation\_bound.by}.
    If a final groupby is present, these labels are stored in its materialized output columns.
    They must not be obtained by re-evaluating the original expressions after the aggregation.

    The preconditions of \ref{line:truncate-id-bound} are satisfied on line \ref{line:id-bound-by}
    (\texttt{id\_bound.by} is equal to \texttt{truncation.by}),
    so by the postcondition of \texttt{truncate\_id\_bound},
    \begin{align*}
        ||d_{\mathrm{Sym}}(\texttt{function}(x)_{g}, \texttt{function}(x')_{g})||_\infty &\le \texttt{row\_bound.per\_group}, \\
        ||d_{\mathrm{Sym}}(\texttt{function}(x)_{g}, \texttt{function}(x')_{g})||_0 &\le \texttt{row\_bound.num\_groups},
    \end{align*}
    where $\texttt{row\_bound}$ denotes the return value.

    Input bound lookup on line \ref{line:id-bound-by} uses the certified execution expressions,
    so that \texttt{truncate\_id\_bound} receives matching input grouping keys.
    When a final groupby is present, line \ref{line:output-bound-by} replaces each key by \texttt{col(output\_name)};
    for filter-only truncations, the keys are unchanged.
    The matcher permits only filters followed by an optional final groupby.
    The constructor checks each preceding filter's keys against the final groupby's keys.
    Thus all translated keys refer to values materialized by that same final groupby.
    This preserves the numerical bounds while expressing them for the output rows.

    For each truncation on line \ref{line:truncation-bound},
    \rustdoc{transformations/make_stable_lazyframe/truncate/fn}{truncate\_id\_bound} on line \ref{line:truncate-id-bound}
    computes upper bounds on the resulting distance between adjacent datasets.
    All acquired bounds are valid upper bounds on the distance between the two datasets,
    since filters only remove rows and the subset check on line \ref{line:allowed_keys_check}
    ensures that a final group-by retains the grouping labels of the preceding filters.

    It is shown that for every pair of elements $x, x'$ in \texttt{input\_domain} and for every pair $(\din, \dout)$,
    where \din\ has the associated type for \texttt{input\_metric} and \dout\ has the associated type for \\ \texttt{output\_metric},
    if $x, x'$ are \din-close under \texttt{input\_metric}, $\texttt{stability\_map}(\din)$ does not raise an exception,
    and $\texttt{stability\_map}(\din) \leq \dout$,
    then $\function(x), \function(x')$ are $\dout$-close under \texttt{output\_metric}.
\end{proof}
\end{document}
